\section{Introduction}\label{sec:intro}

Controlling object trajectories in video generation~\cite{Motionctrl, DragNUWA, DragAnything, TrackGo} is a fundamental task with wide-ranging applications in computer graphics, virtual reality, and interactive media. Precise trajectory control allows for generation of dynamic scenes where objects move in desired paths, enabling creators to create realistic and compelling visual content. Such control is crucial for tasks like animating characters in a virtual environment, simulating physical phenomena, and developing advanced visual effects that require objects to interact seamlessly within a scene.

Despite its importance, controlling object trajectories in video synthesis presents significant challenges. Traditional methods~\cite{DragNUWA, DragAnything, TrackGo} often rely on 2D trajectory inputs drawn directly on images. While these approaches allow for motion representation to some extent, they inherently suffer from ambiguity and complexities associated with interpreting 2D motions in 3D space. Consider the example of animating a hot air balloon flowing over a building as illustrated in \cref{fig:teaser}. A 2D trajectory drawn on the image cannot distinguish whether the balloon should pass in front of or behind the building. This ambiguity arises because a single 2D path can correspond to multiple 3D trajectories due to the lack of 3D information, making it insufficient for precise control over object movements in a 3D space. However, extracting accurate 3D trajectories poses additional difficulties, especially in scenes with occlusions or complex interactions between objects.  For users, inputting valid 3D trajectories is also non-trivial. It often demands specialized knowledge and tools to define object paths accurately within a 3D space, which can be a barrier for artists and non-expert users aiming to create video content.

To address these challenges, we propose \method, a novel model that fine-tunes pre-trained video generation models to incorporate an efficient and effective 3D trajectory control mechanism. 
Our approach introduces an innovative representation of control signal by combining depth information with K-means clustered points of object masks in video. Such control signal can clearly indicate the occlusion and depth changes between objects through the aggregation or separation of clustered points and their depth.
This fusion also captures essential 3D attributes of objects' trajectory without the need for explicit 3D trajectory estimation, thus simplifying the modeling of complex object motions and interactions. For training, we utilize the recently released high-quality Video Object Segmentation (VOS) dataset from SAM2~\cite{SAM2}, which provides rich annotations conducive to our method. By integrating depth cues with clustered points, our representation effectively encodes the object's spatial movements and depth variations over time. This method not only enhances the model's ability to interpret and generate accurate 3D motions but also mitigates issues related to occlusions and depth ambiguities.

We also design a user-friendly inference pipeline that lowers the barrier for users to input 3D trajectories. Users can simply draw trajectories on 2D images and adjust point depths interactively, which the system then interprets as 3D paths for object movements. This approach streamlines the process, making it accessible to users without extensive technical expertise in 3D modeling or animation. 

Our method demonstrates superior performance both quantitatively and qualitatively compared to existing approaches. We achieve accurate 3D trajectory control in image-to-video synthesis task where previous baselines fail.
In summary, our contributions are as follows: We introduce \method, a novel method for controlling 3D object trajectories in video synthesis by combining depth information with K-means clustered points without the need for explicit 3D trajectory tracking. We leverage the high-quality SA-V dataset for training, effectively capturing complex object motions and interactions in diverse scenes. We develop a user-friendly inference pipeline that simplifies the input of 3D trajectories, making it accessible to a broader range of users. To the best of our knowledge, this work is the first to introduce 3D object trajectory control in image-to-video synthesis, paving the way for more advanced and accessible video generation techniques.